%!TEX root = TQFTbook.tex
%%%%%%%%%%%%%%%%%%%%%%%%%%%%%%
In this part, we start building the theory of extended TQFTs, i.e., theories
where we consider not just cobordisms but manifolds with corners. In this part
we are mostly working with 2-dimensional extended TQFTs, leaving higher
dimensional case for the following part.

To give an accurate definition of such an extended TQFT in dimension 2, one
needs to introduce the language of (weak) 2-categories, which we do in
\chref{c:2categories}. We also introduce the notion of a symmetric monoidal
2-category. Our key examples of 2-categories are the 2-category of small
categories (or a variation of it, the 2-category of all locally finite
$\kk$-linear abelian categories) and the category $\Alg$, whose objects are
algebras over $\kk$ and 1-morphisms are bimodules. Both 2-categories play an
important role in future constructions.

Next, in \chref{c:duality-2cat} we discuss dualities in 2-categories, discussing
both the notion of a dual object (which only makes sense in a monoidal
2-category) and of a dual 1-morphism; in particular, we show that for the
2-category of small categories, where 1-morphisms are functors between
categories, so defined duality coincides with the classical notion of
adjointness for functors.
In \chref{c:separable-frobenius}, we study the duality in the 2-category $\Alg$
and show that fully dualizable objects in $\Alg$ are the separable algebras;
moreover, fixing a duality between $\ev$ and $\coev$ 1-morphisms is equivalent
to choosing a structure of a Frobenius algebra, which provides the link to the
construction of (non-extended) TQFTs done earlier.


After completing all these preliminaries, we finally define the notion of an
extended 2d TQFT in \chref{c:2d-extended}. Following work of Schommer--Pries, we
construct the set of generators and relations for the 2-category $\Bord_2$ of
2-dimensional bordisms and use it to show that extended 2d TQFTs are classified
by separable symmetric Frobenius algebras.

Finally, in \chref{c:PL}--\chref{c:2d-lattice} we present an alternative
approach to construction of extended 2d TQFTs. Instead of using Morse functions
to create a presentation of a cobordism as a composition of elementary
cobordisms, we use triangulations (or more generally, PL cell decompositions);
instead of Cerf moves, we use Pachner moves. We show how in dimension 2, it
gives a very explicit construction of an extended 2d TQFT from a separable
symmetric Frobenius algebra $A$, providing a different way of proving the same
classification result.